% methodology_suite.tex -- methods for the standardized 2v2/4v4/6v6 suite.
% Implementation facts trace to CROSS_SCALE_CANONICAL_RECIPE_V1 and the modules named in comments.

\section{Methodology}
\label{sec:method-suite}

\paragraph{Live dual-branch framing (locked).}
The paper's main empirical claim uses one method at three team sizes: the
dual-branch construction in which Poles~A and~B are treated symmetrically
(same recipe, budget, role rule \(k=\lceil N/3\rceil\), and crossover test).
Within each pole, ATTACK and DEFEND are jointly trained under that shared
procedure; both branches receive PPO, and DEFEND alone has an additional
temporary supervised scaffold. Fully Shared+\(z+r\) (\(\pi(a\mid o,z,r)\))
distilled from those teachers is the shared latent-strategy realization of
that framework when Stage~4 is scored on the dual-branch teachers' own
distinction. Share-Encoder remains a comparison; Role-only \(\pi(a\mid o,r)\)
is the critical ablation of Fully Shared+\(z+r\). Historical asymmetric /
frozen-ATTACK constructions, specialist candidate-selection (SQ), and
shared-\(k\) menus on pre-existing specialists (SCK) are excluded from the
main claim.
% Locked method: DUAL_BRANCH_ROLE_COMPOSITE_V1; k=ceil(N/3)

\subsection{Problem: one team, two opponent regimes}
\label{sec:method-problem}

A team of \(N\in\{2,4,6\}\) agents plays capture-the-flag against opponents
drawn from two regimes, Pole~A and Pole~B, which reward different team
behavior. The goal is a team that holds two \emph{distinct} strategies,
selected by a strategy label \(z\in\{A,B\}\), such that each strategy is the
better choice against its own regime. If the two strategies are not distinct,
choosing \(z\) has no effect, and opponent-conditioned play is impossible
regardless of how many games the team wins.

\paragraph{Why win rate alone is not the target.}
A single dominant policy can win often against both regimes. It then offers
one behavior only: there is nothing to select, and nothing to fall back on
when the opponent changes. We therefore measure whether the strategy choice
\emph{matters}, using the crossover effects
\[
\Delta_A = V(\pi_A,A)-V(\pi_B,A),\qquad
\Delta_B = V(\pi_B,B)-V(\pi_A,B),
\]
where \(V(\pi,j)\) is the win rate of strategy \(\pi\) against regime \(j\).
\(\Delta_A>0\) and \(\Delta_B>0\) together mean that each strategy beats the
other on its own regime, i.e.\ the two strategies are complementary. A
negative \(\Delta_B\) means the Pole-A strategy is at least as good on Pole~B
as the strategy built for Pole~B, so one policy dominates and specialization
has failed. We also report absolute win rates, so that any cost of
specialization in raw task success is visible.

\subsection{Compared systems}
\label{sec:method-systems}

Each baseline is the design a practitioner would most likely reach for when
building an opponent-conditioned team. All systems at a given team size use
the same simulator, the same two regimes, the same teacher checkpoints where
applicable, and the same evaluation seeds.

\begin{itemize}
  \item \textbf{Generalist} (\(\pi_G\)): a single unconditioned policy obtained
  by distilling both repaired specialists into one network with no strategy
  variable \(z\). It uses the same suite dataset and teacher policies as the
  strategy-conditioned sharing baselines, isolating the effect of explicit
  strategy conditioning from dataset and teacher choice.
  \item \textbf{Specialists (no roles)}: one policy trained per regime,
  switched by \(z\). This is the simplest way to obtain two strategies. It
  tests whether training against different opponents is, by itself, enough to
  produce complementary behavior.
  \item \textbf{Share-Encoder} and \textbf{Fully Shared+\(z\)}: the same
  specialists compressed by distillation into one network, either sharing the
  observation encoder (with private strategy bodies) or sharing all weights
  and receiving \(z\) as an input. These arms test whether specialization
  survives that compression when strategy identity remains available.
  \item \textbf{Ours (dual-branch)}: Poles~A and~B are treated symmetrically.
  A deterministic rule, CLOSEST\_DEFENDS(\(k\)), assigns the \(k\) living
  agents nearest their own flag to DEFEND for the whole episode and the
  remaining agents to ATTACK, with \(k=\lceil N/3\rceil\)
  (\(k=1,2,2\) at \(N=2,4,6\)). ATTACK and DEFEND are jointly trained under
  that shared organization (Section~\ref{sec:method-ppo}): both branches
  receive on-policy PPO; DEFEND additionally receives a temporary supervised
  scaffold. Leaving ATTACK frozen is an ablation, not this method.
\end{itemize}

\begin{table}[t]
\centering
\caption{Distillation-family baselines. All three students share teachers
\(\pi_A,\pi_B\) and the suite dataset; only architecture and whether \(z\) is
provided differ.}
\label{tab:distill-family}
\begin{tabular}{lccc}
\hline
System & Training source & Uses \(z\)? & Role allocation at test? \\
\hline
Generalist & Distilled from \(\pi_A,\pi_B\) & No & No \\
Share-Encoder & Distilled from \(\pi_A,\pi_B\) & Yes & No \\
Fully Shared+\(z\) & Distilled from \(\pi_A,\pi_B\) & Yes & No \\
\hline
\end{tabular}
\end{table}

The comparisons are designed so that each one changes a single factor:
Specialists versus Ours isolates role organization; Specialists versus the
sharing arms isolates parameter sharing; Generalist versus Fully Shared+\(z\)
isolates explicit strategy conditioning under matched teachers and data.

\subsection{Training the specialists}
\label{sec:method-ppo}

\paragraph{PPO specialist training.}
Each opponent-specific specialist is trained with centralized training and
decentralized execution (CTDE) using an in-house PPO multi-agent learner
(\texttt{rl/custom\_ppo/}). Blue agents share a single actor; a centralized
value function is used only during training. The actor receives a
\(7{\times}20{\times}20\) semantic observation grid, a \(20\)-dimensional
scalar feature vector, and action-validity masks, and outputs a factorized
\(\mathrm{MultiDiscrete}([5,50])\) action (macro-action and waypoint). The
actor is a three-layer CNN (channels \(32\!\to\!64\!\to\!64\)) followed by a
\(256\)--\(256\) MLP; the critic maps a \(34\)-dimensional global state to
\(V(s)\). Training uses generalized advantage estimation with
\(\gamma=0.995\) and \(\lambda=0.99\), together with a shared shaped team
reward that combines terminal outcome (\(+1/-1/-0.5\) for win/lose/draw),
dense task progress (weight \(0.25\); PBRS coefficients \(0.5\) with
PBRS \(\gamma=0.995\); small team presence terms \(0.03/0.02/0.02\)),
event rewards (flag pickup \(0.1\), carry-home \(0.5\); sparse capture
\(\pm100\) points scaled by \(1/100\); tag/kill and opponent-OOB event
rewards zeroed for specialists), and penalties (failed commit \(-0.004\),
own OOB \(-70\) points, stalemate \(-0.08\), spin/idle coefficients
\(0.05/0.03\)), followed by \(\tanh(\cdot/4)\) and clipping to
\([-1,1]\). Strategy-anchored SA-PPO is disabled for every suite teacher.
Each base specialist is trained exclusively against its associated opponent
regime for approximately \(10^6\) environment steps.

The same lineage is applied at every team size; only \(N\), \(k\), seeds,
and tensor sizes that follow from \(N\) differ:
\begin{enumerate}
  \item \textbf{Base specialist} (1\,M steps) against its regime.
  \item \textbf{Entity-repair continuation}: an additive teammate/opponent
  geometry encoding \(g(T,E)\) is added to the actor input
  (\texttt{entity\_residual.py}). It is bias-free, so
  \(g(\emptyset,\emptyset)=0\) exactly and the original observation pathway
  is unchanged.
  \item \textbf{Dual-branch stage} (200\,k joint steps, applied symmetrically
  to A and to B): both \(\pi^{\mathrm{ATTACK}}_s\) and
  \(\pi^{\mathrm{DEFEND}}_s\) are warm-started from the same foundation
  specialist and updated in one shared rollout under CLOSEST\_DEFENDS.
  ATTACK-assigned ticks update only the ATTACK network with standard PPO;
  DEFEND-assigned ticks update only the DEFEND network with the same PPO
  plus a temporary N-prime supervised scaffold (cross-entropy on
  macro-action and waypoint, weight annealed \(0.1\to0\) over
  50\,k--150\,k steps; \texttt{defend\_teacher.py}). ATTACK is not frozen;
  the scaffold is DEFEND-only and temporary.
\end{enumerate}

\subsection{Building the distillation-family baselines}
\label{sec:method-distill}

To compare strategy conditioning and parameter sharing without re-tuning
reinforcement learning per arm, the Generalist, Share-Encoder, and Fully
Shared+\(z\) students are fit by distillation from the same frozen teachers on
the same frozen dataset (Table~\ref{tab:distill-family}). Per team size, one
dataset of decision-time observations (96 episodes per regime) is collected
with CLOSEST\_DEFENDS active, storing entity tensors and roles
(\texttt{collect\_suite\_distillation\_states.py}). Strategy-conditioned
students minimize
\[
\mathcal{L}
=\tfrac12\,\mathrm{KL}\!\big(\pi_A(\cdot\mid o)\,\|\,\pi(\cdot\mid o,z_A)\big)
+\tfrac12\,\mathrm{KL}\!\big(\pi_B(\cdot\mid o)\,\|\,\pi(\cdot\mid o,z_B)\big)
\]
on decision-eligible action heads; the Generalist (and Role-only) use the same
mixture of targets but with a single unconditioned student \(\pi(\cdot\mid o)\)
(no \(z\)). For every retained observation, a strategy-conditioned student is
evaluated under both \(z_0\) and \(z_1\), so parameters shared across modes
receive gradients from both teacher terms, while any mode-specific parameters
are updated only through their corresponding term. Student actors are
initialized independently of the teachers (fresh weights; architecture copied
from the Pole-A teacher), so transfer must arise from the distillation signal
rather than weight inheritance. The optimization recipe is identical for every
arm: Adam with learning rate \(3\times10^{-4}\) and weight decay \(0\),
minibatch size \(256\) stratified as half Pole-A and half Pole-B origin rows
per update (permutation seeded by the arm's training seed), \(20\) epochs,
gradient-norm clip \(1\) on actor parameters only (critic frozen), and terminal
weights kept with no checkpoint selection. At evaluation,
Share-Encoder and Fully Shared+\(z\) receive a forced \(z\); the role rule is
not re-applied, so \(z\) is their only strategy signal. The Generalist has no
\(z\) pathway.

\subsection{Evaluation protocol}
\label{sec:method-eval}

Each system plays every (strategy, regime) cell on the same block of 128
seeds per team size, giving paired per-seed win differences in
\(\{-1,0,+1\}\). Tables report mean \(\pm\) sample standard deviation across
seeds; figures report paired bootstrap 95\% intervals
(\(20{,}000\) resamples). Comparisons labeled ``matched'' use identical seed
blocks. Regime identities are verified at run time, and every run is sealed
against a frozen specification, so no post-hoc seed or checkpoint selection
is possible for the primary \(n=128\) results. Deployment noise (localization
error, motion error, and two-decision control delay) is applied only at
evaluation, on the final system, with a same-block nominal run for paired
comparison.

\paragraph{Post-hoc top-50 subset (descriptive only).}
As a secondary view, seeds may be ranked by deployed score
\(\mathrm{outcome}(A{\to}A)+\mathrm{outcome}(B{\to}B)\) (ties by ascending
seed ID) and the top 50 retained per method and team size. Noise uses the
nominal top-50 seed IDs only, carried into the perturbed conditions so
comparisons stay paired. This subset does not replace the sealed \(n=128\)
tables; regenerators are \texttt{build\_true\_top50.py} (tables and plots
\texttt{fig\_*\_true\_top50}).
Bootstrap intervals on a post-hoc ranked subset are descriptive summaries of
that subset; they do not account for the selection procedure and must not be
labeled as statistically supported confirmatory effects.
